%==============================================================================
% Certificats de primalité inconditionnels pour p = 3m(m+1)+1
% H. Bakkaoui — 2026
%==============================================================================
\documentclass[11pt,a4paper]{article}

\usepackage[utf8]{inputenc}
\usepackage[T1]{fontenc}
\usepackage[french]{babel}
\usepackage{amsmath,amssymb,amsthm,mathtools}
\usepackage{geometry}
\geometry{margin=2.5cm}
\usepackage{graphicx}
\usepackage{booktabs}
\usepackage{array}
\usepackage{enumitem}
\usepackage{caption}
\usepackage[dvipsnames]{xcolor}
\usepackage{tcolorbox}
\tcbuselibrary{skins,breakable}
\usepackage{hyperref}
\hypersetup{colorlinks=true,linkcolor=MidnightBlue,
            citecolor=MidnightBlue,urlcolor=MidnightBlue}
\usepackage{cleveref}

%--- Environnements théorème --------------------------------------------------
\newtheorem{theoreme}{Théorème}[section]
\newtheorem{proposition}[theoreme]{Proposition}
\newtheorem{lemme}[theoreme]{Lemme}
\newtheorem{corollaire}[theoreme]{Corollaire}
\theoremstyle{definition}
\newtheorem{definition}[theoreme]{Définition}
\newtheorem{probleme}[theoreme]{Problème ouvert}
\theoremstyle{remark}
\newtheorem{remarque}[theoreme]{Remarque}

%--- Boîtes colorées ----------------------------------------------------------
\newtcolorbox{rigorousbox}[1][]{
  colback=Green!8, colframe=Green!60!black, arc=1mm,
  title=\textbf{Résultat rigoureux}, fonttitle=\bfseries, breakable, #1}
\newtcolorbox{resultbox}[1][]{
  colback=RoyalBlue!5, colframe=RoyalBlue!70!black, arc=1mm,
  title=\textbf{Résultat calculatoire}, fonttitle=\bfseries, breakable, #1}
\newtcolorbox{openprobbox}[1][]{
  colback=Goldenrod!10, colframe=Goldenrod!70!black, arc=1mm,
  title=\textbf{Problème ouvert}, fonttitle=\bfseries, breakable, #1}

%--- Raccourcis ---------------------------------------------------------------
\newcommand{\N}{\mathbb{N}}
\newcommand{\Z}{\mathbb{Z}}
\newcommand{\F}{\mathcal{F}}
\newcommand{\Leg}[2]{\left(\dfrac{#1}{#2}\right)}

%==============================================================================
\title{Certificats de primalité inconditionnels\\
pour les entiers de la forme $p = 3m(m+1)+1$ :\\[4pt]
{\large Filtres arithmétiques, recherche multi-cœurs\\
et records jusqu'à $29\,998$ chiffres décimaux}}

\author{Hassane \textsc{Bakkaoui}\\
\small Chercheur indépendant\\
\small \texttt{bakkahassa@hotmail.com}}

\date{2026}
%==============================================================================

\begin{document}
\maketitle

%------------------------------------------------------------------------------
\begin{abstract}
Nous étudions la sous-famille paramétrique $p = 3m(m+1)+1$,
$m = 2^a\cdot 3^b - 1$, $a,b \in \N^*$, du point de vue de la
certification \emph{inconditionnelle} de primalité via le critère de
Pocklington--Lehmer \cite{Pocklington1914,LehmerPSW1975}.

La structure $3$-lisse de $m+1 = 2^a\cdot 3^b$ fournit, pour tout $(a,b)$,
un diviseur complètement factorisé $F = 2^a\cdot 3^{b+1}$ de $p-1$
satisfaisant $F>\sqrt{p}$ inconditionnellement
(Lemme~\ref{lem:factorisation}), réduisant le certificat à la recherche
de témoins pour $q\in\{2,3\}$ seulement.

Nous établissons trois théorèmes de structure arithmétique :
\begin{itemize}
\item \textbf{Théorème~\ref{thm:mod6} :}
  $p\equiv 1\pmod{6}$ pour tout $(a,b)$.
\item \textbf{Théorème~\ref{thm:cubique} :}
  tout premier $q\equiv 2\pmod{3}$ vérifie $q\nmid p$,
  réduisant le crible de $50\%$.
  Preuve : $-3$ est un non-résidu quadratique mod $q$ lorsque
  $q\equiv 2\pmod{3}$.
\item \textbf{Théorème~\ref{thm:mod7} :}
  $7\mid p$ si et seulement si $(a\bmod 3,\,b\bmod 6)\in\F$,
  où $\F$ est un ensemble explicite de \emph{six} paires,
  éliminant $33\%$ des candidats à coût arithmétique nul.
  Preuve : $\mathrm{ord}_7(2)=3$, $\mathrm{ord}_7(3)=6$.
\end{itemize}

Combinés, ces trois résultats éliminent $\approx 87\%$ de tous les
candidats avant tout calcul en grande précision. L'algorithme
multi-cœurs \textsc{PrimeQuest} (versions v1--v5) exploite ce pipeline
pour produire quatre certificats de Pocklington--Lehmer inconditionnels,
le plus grand étant
\[
  p = 3m(m+1)+1,\quad m = 2^{19435}\cdot 3^{19173}-1,
\]
un nombre premier de $\mathbf{29\,998}$ \textbf{chiffres décimaux} trouvé
en $3$ h $06$ min sur un ordinateur portable grand public
(ARM64, 9 cœurs).

\smallskip
\noindent\textbf{MSC 2020 :} 11A41, 11Y11, 11N13.\\
\textbf{Mots-clés :} certificat de primalité, Pocklington--Lehmer,
premiers hexagonaux centrés, filtres arithmétiques, résidus quadratiques,
records de grands premiers.
\end{abstract}

\tableofcontents

%==============================================================================
\section{Introduction}\label{sec:intro}
%==============================================================================

Prouver inconditionnellement qu'un entier donné est premier est une tâche
centrale de la théorie algorithmique des nombres. Pour les grands premiers
génériques, la voie la plus pratique est le critère de
Pocklington--Lehmer~\cite{Pocklington1914,LehmerPSW1975} : exhiber un
diviseur complètement factorisé $F>\sqrt{p}$ de $p-1$ et vérifier des
témoins. La difficulté principale réside dans la construction d'un tel $F$.

La sous-famille
\[
  p = 3m(m+1)+1,\qquad m = 2^a\cdot 3^b-1,\qquad a,b\in\N^*,
\]
est construite de sorte que $m+1=2^a\cdot 3^b$ soit $3$-lisse par
construction, rendant $F=2^a\cdot 3^{b+1}$ complètement factorisé et
vérifiant $F>\sqrt{p}$ pour tout $(a,b)$, sans aucune condition sur $p$.

Au-delà de la structure du certificat, la famille satisfait trois
contraintes arithmétiques indépendantes
(Théorèmes~\ref{thm:mod6}--\ref{thm:mod7}) qui éliminent collectivement
$87\%$ des candidats avant tout calcul en grande précision. En exploitant
ces trois résultats, l'algorithme \textsc{PrimeQuest} a produit quatre
certificats inconditionnels jusqu'à $29\,998$ chiffres décimaux.

\smallskip
\noindent\textbf{Lien avec les articles compagnons.}
Cet article est auto-contenu. Les aspects distributionnels et analytiques
de la famille (loi des gaps, constante de Hardy--Littlewood, connexion à
$\zeta(s)$) sont traités dans~\cite{PaperI,PaperII}.

%==============================================================================
\section{La sous-famille de Pocklington}\label{sec:famille}
%==============================================================================

\subsection{Définition et nombre de chiffres}

\begin{definition}[Sous-famille de Pocklington]\label{def:pocklington}
Pour des entiers $a,b\geq 1$, on définit
\[
  m(a,b) = 2^a\cdot 3^b - 1, \qquad
  p(a,b) = 3m(m+1)+1 = 3m^2+3m+1.
\]
Le nombre de chiffres décimaux de $p$ est
$D(a,b) = \lfloor 2a\log_{10}2 + (2b+1)\log_{10}3\rfloor + 1.$
\end{definition}

\subsection{Factorisation de $p-1$ et structure du certificat}

\begin{rigorousbox}
\begin{lemme}[Factorisation explicite]\label{lem:factorisation}
Pour $p$ comme dans la Définition~\ref{def:pocklington},
\[
  p - 1 = 2^a\cdot 3^{b+1}\cdot m, \qquad F := 2^a\cdot 3^{b+1} > \sqrt{p}.
\]
\end{lemme}
\end{rigorousbox}

\begin{proof}
Comme $m+1=2^a\cdot 3^b$,
\[
  p-1 = 3m(m+1) = 3m\cdot 2^a\cdot 3^b = 2^a\cdot 3^{b+1}\cdot m.
\]
Comme $p = 3m^2+3m+1 < 3(m+1)^2$, on obtient $\sqrt{p}<\sqrt{3}(m+1)
<3(m+1)= 3\cdot 2^a\cdot 3^b = 2^a\cdot 3^{b+1} = F$.
\end{proof}

\begin{rigorousbox}
\begin{theoreme}[Certificat de Pocklington--Lehmer
\cite{Pocklington1914,LehmerPSW1975}]\label{thm:Pocklington}
Soit $p = p(a,b)$ comme dans la Définition~\ref{def:pocklington}.
Alors $p$ est premier si et seulement s'il existe des entiers $w_2,w_3$
tels que, pour $q\in\{2,3\}$,
\[
  w_q^{p-1}\equiv 1\pmod{p}\quad\text{et}\quad
  \gcd\!\bigl(w_q^{(p-1)/q}-1,\,p\bigr)=1.
\]
\end{theoreme}
\end{rigorousbox}

\begin{corollaire}\label{cor:pocklington}
Le Théorème~\ref{thm:Pocklington} s'applique inconditionnellement à tout
candidat de la sous-famille : par le Lemme~\ref{lem:factorisation},
$F>\sqrt{p}$ est toujours vérifié, et le certificat requiert au plus
deux témoins ($q\in\{2,3\}$ seulement).
\end{corollaire}

\begin{remarque}
Dans les quatre certificats produits (Section~\ref{sec:resultats}),
les témoins $w_2 = 5$ et $w_3 = 7$ ont réussi sans exception.
\end{remarque}

%==============================================================================
\section{Théorèmes de structure arithmétique}\label{sec:filtres}
%==============================================================================

Nous établissons trois contraintes indépendantes sur la famille
$p = 3m(m+1)+1$ avec $m = 2^a\cdot 3^b-1$. Chacune est prouvée
rigoureusement et traduite en économie calculatoire concrète.

\subsection{Congruence modulo 6}

\begin{rigorousbox}
\begin{theoreme}[Modulo $6$]\label{thm:mod6}
Pour tout $a,b\geq 1$, $p(a,b)\equiv 1\pmod{6}$.
\end{theoreme}
\end{rigorousbox}

\begin{proof}
\textit{Modulo $2$.} Comme $m = 2^a\cdot 3^b-1$ est impair, $m(m+1)$
est pair, donc $3m(m+1)$ est pair et $p = 3m(m+1)+1$ est impair.

\textit{Modulo $3$.} On a $3\mid 3m(m+1)$, donc $p\equiv 1\pmod{3}$.

Ensemble, $p\equiv 1\pmod{2}$ et $p\equiv 1\pmod{3}$ donnent
$p\equiv 1\pmod{6}$ par le théorème chinois des restes.
\end{proof}

\begin{corollaire}
Aucun premier $q\in\{2,3\}$ ne divise $p(a,b)$, et $p$ appartient à la
classe admissible $+1$ modulo $6$ pour tout $(a,b)$.
\end{corollaire}

\subsection{Crible par réciprocité cubique}

\begin{rigorousbox}
\begin{theoreme}[Crible par réciprocité cubique]\label{thm:cubique}
Soit $m\in\N^*$, $p = 3m(m+1)+1$, et $q\geq 5$ un premier avec
$q\equiv 2\pmod{3}$. Alors $q\nmid p$.
\end{theoreme}
\end{rigorousbox}

\begin{proof}
La condition $q\mid p$ est équivalente à $3m^2+3m+1\equiv 0\pmod{q}$.
En complétant le carré, on obtient $(6m+3)^2\equiv -3\pmod{q}$, donc
$-3$ doit être un résidu quadratique modulo $q$. Par la multiplicativité
du symbole de Legendre et la réciprocité quadratique~\cite{IrelandRosen},
\[
  \Leg{-3}{q} = \Leg{-1}{q}\Leg{3}{q}
  = (-1)^{\frac{q-1}{2}}\cdot\Leg{q}{3}(-1)^{\frac{q-1}{2}\cdot\frac{3-1}{2}}.
\]
Un calcul direct montre $\bigl(\tfrac{-3}{q}\bigr)=+1$ si $q\equiv 1\pmod{3}$
et $\bigl(\tfrac{-3}{q}\bigr)=-1$ si $q\equiv 2\pmod{3}$. Dans ce dernier
cas, $3m^2+3m+1\equiv 0\pmod{q}$ n'a pas de solution, donc $q\nmid p$.
\end{proof}

\begin{corollaire}[Réduction du crible de $50\%$]\label{cor:crible-50}
Dans la division d'essai jusqu'à une borne $L$, seuls les premiers
$q\equiv 1\pmod{3}$ doivent être testés. Par le théorème de Dirichlet,
ceux-ci représentent la moitié de tous les premiers dans $[5,L]$.
Concrètement : des $78\,498$ premiers dans $[5,10^6]$, seulement
$39\,231$ vérifient $q\equiv 1\pmod{3}$.

\smallskip
\noindent\emph{Note.} Le Théorème~\ref{thm:cubique} s'applique à la
famille générale $p=3m(m+1)+1$ avec $m\in\N^*$ arbitraire,
indépendamment de la structure $3$-lisse.
\end{corollaire}

\subsection{Classes interdites modulo $7$}

\begin{rigorousbox}
\begin{theoreme}[Classes interdites modulo $7$]\label{thm:mod7}
Soit $m = 2^a\cdot 3^b-1$, $a,b\geq 1$. On définit
\[
  \F_7 = \bigl\{(0,2),(0,3),(1,0),(1,1),(2,4),(2,5)\bigr\}
  \;\subset\; \{0,1,2\}\times\{0,1,2,3,4,5\}.
\]
Alors
\[
  7\mid p(a,b)\quad\Longleftrightarrow\quad
  (a\bmod 3,\;b\bmod 6)\in\F_7.
\]
En particulier, exactement $|\F_7|/18 = 1/3$ de toutes les classes
sont interdites.
\end{theoreme}
\end{rigorousbox}

\begin{proof}
Comme $\mathrm{ord}_7(2)=3$ et $\mathrm{ord}_7(3)=6$, la valeur
$m+1 = 2^a\cdot 3^b \bmod 7$ ne dépend que de $r=a\bmod 3$ et
$s=b\bmod 6$. En posant $t = t(r,s):= 2^r\cdot 3^s\bmod 7$,
on a $m\equiv t-1\pmod{7}$ et
\[
  p \equiv 3(t-1)t+1 = 3t^2-3t+1\pmod{7}.
\]
La condition $7\mid p$ devient $3t^2-3t+1\equiv 0\pmod{7}$.
En multipliant par $5\equiv 3^{-1}\pmod{7}$ :
\[
  t^2-t-2\equiv 0\pmod{7},\qquad(t-2)(t+1)\equiv 0\pmod{7},
\]
donc $t\equiv 2$ ou $t\equiv 6\pmod{7}$. Le Tableau~\ref{tab:valeurs-t}
liste les $3\times 6=18$ valeurs de $t(r,s)$ ; exactement les six paires
de $\F_7$ donnent $t\in\{2,6\}$.
\end{proof}

\begin{table}[ht]\centering
\caption{Valeurs de $t(r,s)=2^r\cdot 3^s\bmod 7$.
Les entrées dans $\{2,6\}$ (en gras) correspondent à $7\mid p$ ;
ce sont les six paires de $\F_7$.}
\label{tab:valeurs-t}
\begin{tabular}{@{}c|cccccc@{}}
\toprule
$r\;\backslash\; s$ & $0$ & $1$ & $2$ & $3$ & $4$ & $5$ \\
\midrule
$0$ & $1$ & $3$ & $\mathbf{2}$ & $\mathbf{6}$ & $4$ & $5$ \\
$1$ & $\mathbf{2}$ & $\mathbf{6}$ & $4$ & $5$ & $1$ & $3$ \\
$2$ & $4$ & $5$ & $1$ & $3$ & $\mathbf{2}$ & $\mathbf{6}$ \\
\bottomrule
\end{tabular}
\end{table}

\begin{corollaire}[Élimination gratuite de $33\%$]\label{cor:t2-gratuit}
Vérifier $(a\bmod 3,\,b\bmod 6)\in\F_7$ ne requiert que deux réductions
modulaires et une recherche dans un ensemble de six paires — aucune
arithmétique sur des entiers à $D$ chiffres. Lorsque la recherche parcourt
les paires $(a,b)$ uniformément sur les $18$ classes de résidus
(ce qui est le cas dans la recherche en zigzag de la
Section~\ref{sec:algo}), exactement $1/3$ de tous les candidats sont
rejetés à coût nul.
\end{corollaire}

\subsection{Taux d'élimination combiné}

\begin{proposition}[Efficacité combinée des filtres]\label{prop:combine}
Parmi tous les couples $(a,b)$ candidats explorés par la recherche en
zigzag :
\begin{enumerate}[label=(\roman*)]
\item Le Théorème~\ref{thm:mod7} élimine $33{,}3\%$ (consultation, coût nul).
\item Le Théorème~\ref{thm:cubique} réduit le crible à $q\equiv 1\pmod{3}$ ;
  parmi les $66{,}7\%$ restants, le crible élimine encore $\approx 80\%$.
\item La fraction atteignant Miller--Rabin est $\approx 13\%$.
\end{enumerate}
Pour le record à $29\,998$ chiffres ($2\,221$ couples testés) :
$33{,}3\%$ par le Théorème~\ref{thm:mod7},
$53{,}8\%$ par le crible,
$\mathbf{87{,}1\%}$ au total.
\end{proposition}

%==============================================================================
\section{L'algorithme \textsc{PrimeQuest}}\label{sec:algo}
%==============================================================================

\subsection{Calibration du centre et zigzag}

Pour un nombre cible de $D$ chiffres, le \emph{centre de recherche} est
\[
  a_0 = b_0 = \Bigl\lfloor\frac{D}{2\log_{10}6}\Bigr\rfloor \approx 0{,}643\,D.
\]
Le zigzag s'étend vers l'extérieur : pour $\delta = 0, 1, 2, \ldots$,
il alterne $a = a_0+\delta$ et $a = a_0-\delta$, en énumérant pour chaque
$a$ tous les $b$ avec $|D(a,b)-D|\leq\tau$ (tolérance $\tau=10$).

\begin{remarque}
Observation empirique : $b/a = 1{,}083$ pour $D=19\,999$ ;
$b/a = 0{,}987$ pour $D=29\,998$. Comme ce ratio varie entre les premiers,
le choix symétrique $a_0=b_0$ minimise la distance dans le pire cas à
la zone des premiers.
\end{remarque}

\subsection{Pipeline à quatre étapes}

Chaque couple $(a,b)$ passe par :
\begin{enumerate}
\item \textbf{Filtre mod-$7$} (Thm~\ref{thm:mod7}) : rejeter si
  $(a\bmod 3,b\bmod 6)\in\F_7$. Coût : $O(1)$.
\item \textbf{Crible} (Cor~\ref{cor:crible-50}) : pour $q\equiv 1\pmod{3}$,
  $q\leq L$, calculer $p\bmod q$ via des exponentiations modulaires
  sur petits entiers ; rejeter si $p\equiv 0$.
\item \textbf{Miller--Rabin} : $k$ tours sur $p(a,b)$ en pleine précision.
\item \textbf{Pocklington} : si MR réussit, chercher des témoins
  $w_2,w_3$ du Théorème~\ref{thm:Pocklington}.
\end{enumerate}

\subsection{Exécution multi-cœurs et point de reprise}

Les couples sont distribués entre les cœurs via un pool à mémoire partagée.
Les workers ignorent SIGINT/SIGTERM ; la sauvegarde de point de reprise
toutes les $60$ s permet un redémarrage exact après interruption.

\subsection{Versions de l'algorithme}

\begin{table}[ht]\centering
\caption{Versions successives de \textsc{PrimeQuest}.}
\label{tab:versions}
\begin{tabular}{@{}clcc@{}}
\toprule
Version & Optimisation ajoutée & Tours MR $k$ & Limite crible $L$ \\
\midrule
v1 & Recherche séquentielle de base & $25$ & $10^6$ \\
v2 & Crible $\to$ $q\equiv 1\pmod{3}$ (Thm~\ref{thm:cubique}) & $25$ & $10^6$ \\
v3 & Multi-cœurs, point de reprise, calibration centre & $25$ & $10^6$ \\
v4 & Filtre mod-$7$ (Thm~\ref{thm:mod7}) & $25$ & $10^6$ \\
v5 & $k\colon 25\to 3+20^{*}$ ; $L\colon 10^6\to 10^7$ & $3+20$ & $10^7$ \\
\bottomrule
\end{tabular}\\[3pt]
\footnotesize $^{*}$ v5 : $3$ tours initiaux, puis $20$ tours de
confirmation avant Pocklington.
\end{table}

\begin{remarque}[Correction de la réduction des tours MR]
Miller--Rabin est \emph{déterministe pour les premiers} : tout nombre
premier passe chaque tour avec certitude. Réduire $k$ de $25$ à $3$ ne
présente aucun risque de manquer un vrai premier ; les composites rares
passant $3$ tours ($\leq 4^{-3}\approx 1{,}6\%$) sont rattrapés par
les $20$ tours de confirmation et le test de Pocklington.
\end{remarque}

%==============================================================================
\section{Résultats calculatoires}\label{sec:resultats}
%==============================================================================

\subsection{Quatre premiers certifiés}

\begin{resultbox}
Le Tableau~\ref{tab:resultats} regroupe les quatre certificats
inconditionnels de Pocklington--Lehmer. Dans chaque cas :
$w_2=5$ pour $q=2$ et $w_3=7$ pour $q=3$.
\end{resultbox}

\begin{table}[ht]\centering
\caption{Certificats Pocklington de la campagne \textsc{PrimeQuest}.}
\label{tab:resultats}
\begin{tabular}{@{}rrrrrrrr@{}}
\toprule
Chiffres de $p$ & $a$ & $b$ & $b/a$ & Tests MR & Élim. & Temps & Workers \\
\midrule
$9\,998$  & $6\,212$  & $6\,738$  & $1{,}085$ & ---      & ---      & ---             & $1$ \\
$10\,000$ & $6\,213$  & $6\,740$  & $1{,}085$ & ---      & ---      & ---             & $1$ \\
$19\,999$ & $12\,228$ & $13\,242$ & $1{,}083$ & $1\,050$ & $87{,}9\%$ & $2$\,h\,$40$\,min & $9$ \\
$29\,998$ & $19\,435$ & $19\,173$ & $0{,}987$ & $286$    & $87{,}1\%$ & $3$\,h\,$06$\,min & $9$ \\
\bottomrule
\end{tabular}
\end{table}

\subsection{Certificat du premier à $29\,998$ chiffres}

\begin{rigorousbox}
Le nombre premier $p = 3m(m+1)+1$ avec $m = 2^{19435}\cdot 3^{19173}-1$
($14\,999$ chiffres) admet le certificat inconditionnel suivant :
\begin{align*}
  p-1 &= F\cdot m,\quad F = 2^{19435}\cdot 3^{19174}
         \;(15\,000\text{ chiffres}),\quad F>\sqrt{p}, \\
  w_2 &= 5:\quad 5^{p-1}\equiv 1\pmod{p},\quad
         \gcd(5^{(p-1)/2}-1,\,p)=1, \\
  w_3 &= 7:\quad 7^{p-1}\equiv 1\pmod{p},\quad
         \gcd(7^{(p-1)/3}-1,\,p)=1.
\end{align*}
Les premiers et derniers chiffres de $p$ sont :
\[
  p = 1602087359509060348500\cdots 16930852697403293697.
\]
\end{rigorousbox}

\subsection{Statistiques des filtres pour le record à $29\,998$ chiffres}

\begin{center}
\begin{tabular}{@{}lrr@{}}
\toprule
Étape & Couples éliminés & Fraction \\
\midrule
Filtre mod-$7$ (Thm~\ref{thm:mod7}) & $740$ & $33{,}3\%$ \\
Crible (Cor~\ref{cor:crible-50}) & $1\,195$ & $53{,}8\%$ \\
Miller--Rabin (composites) & $285$ & $12{,}8\%$ \\
\midrule
\textbf{Total éliminé} & $2\,220$ & $\mathbf{87{,}1\%}$ \\
Premier trouvé & $1$ & $0{,}045\%$ \\
\bottomrule
\end{tabular}
\end{center}

%==============================================================================
\section{Performance et mise à l'échelle}\label{sec:echelle}
%==============================================================================

\subsection{Coût dominant}

Le coût dominant est Miller--Rabin : un tour sur un module à $D$ chiffres
coûte $O(D^2\log D\log\log D)$ opérations binaires (multiplication FFT
de GMP \cite{GMP}). Les vérifications de crible et mod-$7$ sont
négligeables en comparaison.

\subsection{Loi de mise à l'échelle empirique}

Temps d'horloge par événement MR (temps total $\div$ nombre de tests MR),
machine : Dell Inspiron 14 Plus (Snapdragon X Plus ARM64, 9 workers) :

\begin{center}
\begin{tabular}{@{}rrrr@{}}
\toprule
$D$ & Tours $k$ & $\tau_{\text{MR}}$ (horloge) & Version \\
\midrule
$19\,999$ & $25$ & $\approx 9$\,s  & v3 \\
$29\,998$ & $25$ & $\approx 39$\,s & v4 \\
$29\,998$ & $3$  & $\approx 5$\,s  & v5 \\
\bottomrule
\end{tabular}
\end{center}

Le rapport $39/9 \approx 4{,}3$ pour $D = 30\,000/20\,000 = 1{,}5\times$
est cohérent avec la mise à l'échelle empirique $\tau \propto D^{2,5}$.

\subsection{Projection : record à $50\,000$ chiffres}

Avec v5 ($k=3$) et la mise à l'échelle $D^{2,5}$ :
\[
  \tau_{\text{MR}}(50\,000) \approx 5 \times\Bigl(\tfrac{50}{30}\Bigr)^{2,5}
  \approx 18\text{ s/événement}.
\]
Pour $400$--$500$ événements MR attendus : temps total projeté
$\mathbf{6}$--$\mathbf{10}$ \textbf{heures} (9 workers, même machine).
Une recherche est actuellement en cours.

%==============================================================================
\section{Problèmes ouverts}\label{sec:ouverts}
%==============================================================================

\begin{openprobbox}
\begin{probleme}[Filtres arithmétiques généralisés]\label{op:filtre-gen}
Pour les premiers $\ell\equiv 1\pmod{3}$, $\ell\in\{13,19,31,\ldots\}$,
dériver l'ensemble interdit
\[
  \F_\ell = \bigl\{(a\bmod\mathrm{ord}_\ell(2),\;
                    b\bmod\mathrm{ord}_\ell(3)) : \ell\mid p(a,b)\bigr\}.
\]
Pour $\ell=13$ : $\mathrm{ord}_{13}(2)=12$, $\mathrm{ord}_{13}(3)=3$ ;
la période est $36$ et le filtre élimine $2/13\approx 15\%$ de candidats
supplémentaires. La combinaison de filtres pour plusieurs petits $\ell$
pourrait porter le taux d'élimination pré-MR au-dessus de $90\%$.
\end{probleme}
\end{openprobbox}

\begin{openprobbox}
\begin{probleme}[Certificat à $50\,000$ chiffres]\label{op:50k}
Produire un certificat de Pocklington inconditionnel pour un premier de
$50\,000$ chiffres décimaux dans la sous-famille. Par la
Proposition~\ref{prop:combine}, $\approx 87\%$ des candidats sont
pré-filtrés ; le temps projeté est $6$--$10$ heures (v5, ARM64, 9 workers).
\end{probleme}
\end{openprobbox}

\begin{openprobbox}
\begin{probleme}[Universalité des témoins $w_2=5$, $w_3=7$]
\label{op:temoins}
Dans les quatre certificats, $w_2=5$ et $w_3=7$ ont réussi. Prouver
(ou réfuter) que $5$ et $7$ sont des témoins de Pocklington valides
pour \emph{tout} premier de la sous-famille.
\end{probleme}
\end{openprobbox}

%==============================================================================
\section{Conclusion}\label{sec:conclusion}
%==============================================================================

Nous avons donné un traitement auto-contenu des certificats de primalité
inconditionnels pour la sous-famille $3$-lisse $p=3m(m+1)+1$,
$m=2^a\cdot 3^b-1$.

Trois théorèmes arithmétiques élémentaires mais efficaces — l'identité
mod-$6$ (Théorème~\ref{thm:mod6}), le crible par réciprocité cubique
(Théorème~\ref{thm:cubique}), et le filtre des classes interdites mod-$7$
(Théorème~\ref{thm:mod7}) — éliminent ensemble $87\%$ des candidats avant
tout calcul en grande précision. Leurs preuves reposent sur des outils
standard : le théorème chinois des restes, la non-résidualité de $-3$
modulo les premiers $q\equiv 2\pmod{3}$, et les ordres multiplicatifs
$\mathrm{ord}_7(2)=3$, $\mathrm{ord}_7(3)=6$.

Appliqué via l'algorithme multi-cœurs \textsc{PrimeQuest} sur un
ordinateur portable grand public, ce cadre a produit quatre certificats
inconditionnels, culminant en un premier de $\mathbf{29\,998}$ chiffres
décimaux ($286$ événements MR, $3$ h $06$ min). Une cinquième recherche
ciblant $50\,000$ chiffres est en cours.

%==============================================================================
\section*{Remerciements}
%==============================================================================

L'auteur remercie les développeurs de \textsc{gmpy2} et de GMP pour avoir
rendu l'arithmétique en grande précision accessible aux chercheurs
indépendants. Les calculs ont été effectués sur un Dell Inspiron 14 Plus 7441
(Snapdragon X Plus ARM64, 10 cœurs Oryon, mémoire LPDDR5X).

%==============================================================================
\begin{thebibliography}{99}
%==============================================================================

\bibitem{PaperI}
H.~Bakkaoui, \emph{Une famille paramétrique de premiers via les formes
quadratiques : distribution heuristique du décalage minimal, validation
numérique à grande échelle et contrôle statistique des corrélations
spectrales spurieuses},
Document de travail (Article~I), 2026.

\bibitem{PaperII}
H.~Bakkaoui, \emph{Familles de premiers paramétriques via les formes
quadratiques : invariants de couche, preuves conditionnelles et connexions
à la fonction zêta de Riemann},
Document de travail (Article~II), 2026.

\bibitem{AKS2002}
M.~Agrawal, N.~Kayal et N.~Saxena, \emph{PRIMES is in P},
Ann.\ of Math.\ \textbf{160} (2004), 781--793.

\bibitem{GMP}
T.~Granlund et l'équipe de développement GMP,
\emph{GNU Multiple Precision Arithmetic Library}, version~6.3.0,
\texttt{https://gmplib.org}, 2023.

\bibitem{LehmerPSW1975}
J.~Brillhart, D.~H.~Lehmer et J.~L.~Selfridge, \emph{New primality
criteria and factorizations of $2^m\pm 1$},
Math.\ Comp.\ \textbf{29} (1975), 620--647.

\bibitem{Pocklington1914}
H.~C.~Pocklington, \emph{The determination of the prime or composite
nature of large numbers by Fermat's theorem},
Proc.\ Cambridge Philos.\ Soc.\ \textbf{18} (1914--1916), 29--30.

\bibitem{IrelandRosen}
K.~Ireland et M.~Rosen, \emph{A Classical Introduction to Modern
Number Theory}, 2\textsuperscript{e} éd., Springer, 1990.

\end{thebibliography}

\end{document}
